\section{Halation}
\label{sec:halation}

\subsection{Physical Origin}

Light entering the emulsion is not entirely absorbed. Some fraction traverses
all three sensitized layers, enters the transparent base, and is redirected back
into the emulsion by two distinct mechanisms:

\begin{enumerate}
\item \textbf{Base reflection.} At the base--air interface on the rear surface,
rays striking at angles beyond the critical angle
$\theta_c = \arcsin(1/n)$ are totally internally reflected. Having crossed the
base twice, they re-enter the emulsion displaced by
\begin{equation}
r = 2\, t_b \tan\theta,
\label{eq:ringradius}
\end{equation}
where $t_b$ is base thickness. Because only $\theta > \theta_c$ reflects
totally, the returned energy forms an \emph{annulus} of inner radius
$r_{\min} = 2 t_b \tan\theta_c$, not a disc. For a $125\,\mu$m acetate base with
$n = 1.50$, $\theta_c = 41.8^\circ$ and $r_{\min} \approx 224\,\mu$m --- roughly
$0.6\%$ of the width of a 135 frame, which is a visible ring.
\item \textbf{Diffuse scattering.} Light scatters from silver halide crystals
and from layer interfaces, producing a monotone, heavy-tailed spread with no
characteristic ring.
\end{enumerate}

Modern films carry an antihalation undercoat or a rem-jet backing that absorbs
most light reaching the base, suppressing the ring almost entirely and leaving
the diffuse component. \EM{} models both and lets the stock profile weight
them, because the ring is exactly what makes certain older stocks and certain
motion picture negatives look the way they do, while the diffuse component is
what modern stocks exhibit.

\subsection{Wavelength Dependence}

Halation is strongly red-biased, for two compounding reasons. First, the
red-sensitive layer sits deepest in the stack, so returning light passes through
it first and most strongly. Second, longer wavelengths are absorbed less by the
emulsion and by the antihalation layer, so more red light survives the round
trip. The result is that the characteristic halo around a bright highlight is
orange-red, and this is the single most recognizable film cue in the entire
pipeline.

Scattering lengths therefore satisfy $\ell_R > \ell_G > \ell_B$, with typical
ratios $\ell_R : \ell_G : \ell_B \approx 1.00 : 0.62 : 0.44$, and channel
weights $\boldsymbol{\beta} \approx (1.00, 0.42, 0.22)$.

\subsection{The Source Term}

Halation is driven by scene exposure above a threshold, evaluated in the
\emph{linear scene-referred domain} before the characteristic curve, per the
ordering derivation of Section~\ref{sec:arch}. A hard threshold
$\max(0, E - \tau)$ has a discontinuous derivative and produces a visible
contour where the halo begins. \EM{} reuses the softplus primitive of
\eqref{eq:softplus}:

\begin{equation}
\boxed{\;S(\mathbf{p}) \;=\; \left[\softp{\kappa_h}{E_Y(\mathbf{p}) - \tau}\right]^{\,q}\;}
\label{eq:halsource}
\end{equation}

where $E_Y$ is the luminance-weighted exposure, $\tau$ is the threshold in
relative exposure units (default $\tau = 1.6$, i.e. about $3.2$ stops above
middle gray), $\kappa_h$ is the knee softness (default $0.15$), and $q$ is the
source exponent (default $1.0$; $q = 2$ concentrates the effect into the very
brightest sources, which suits point highlights such as street lamps).

Using luminance rather than per-channel exposure is deliberate: the source of the
scattered light is the total radiance at that point, and the channel dependence
enters through the transport, not the source. Per-channel thresholding produces
colored source regions and looks wrong.

\subsection{The Point Spread Function}

The diffuse component is modelled with a normalized isotropic exponential:

\begin{equation}
\boxed{\;h_\ell(r) \;=\; \frac{1}{2\pi \ell^2}\, e^{-r/\ell}, \qquad \int_{\Rspace^2} h_\ell \, dA = 1\;}
\label{eq:halpsf}
\end{equation}

whose two-dimensional Fourier transform is available in closed form,

\begin{equation}
H_\ell(f) \;=\; \left[1 + (2\pi \ell f)^2\right]^{-3/2}.
\label{eq:halotf}
\end{equation}

Equation \eqref{eq:halotf} is what makes the model verifiable: it is the optical
transfer function of the scattering layer, and it can be compared directly
against a measured MTF degradation or against the radial profile of a
photographed point source (AC-4).

The exponential is chosen over a Gaussian because it is heavy-tailed. A Gaussian
halo has a definite edge; an exponential halo has a bright core and a long,
faint skirt, which is what the eye reads as "glow inside the film" rather than
"blur applied to the highlights." This distinction is perceptually decisive and
costs nothing to get right.

The ring component uses

\begin{equation}
h^{\mathrm{ring}}(r) \;=\; \frac{1}{Z}\exp\!\left[-\frac{(r - r_{\min})^2}{2 w_r^2}\right],
\label{eq:halring}
\end{equation}

with $r_{\min}$ from \eqref{eq:ringradius}, $w_r$ the ring width set by the
angular spread, and $Z$ the normalization $\int h^{\mathrm{ring}} dA = 1$. The
composite PSF is

\begin{equation}
h(r) = (1 - \omega)\, h_\ell(r) + \omega\, h^{\mathrm{ring}}(r),
\end{equation}

with $\omega$ the base-reflection weight, near zero for modern stocks and up to
$0.35$ for stocks without effective antihalation.

\subsection{Recombination and Energy Conservation}

The scattered light must be added to the direct exposure --- and, strictly,
removed from it, since it is the same photons. The energy-conserving form is

\begin{equation}
\boxed{\;E'_c \;=\; \left(1 - \alpha_h \beta_c\right) E_c \;+\; \alpha_h \beta_c \left(h_{\ell_c} * S\right)\;}
\label{eq:haladd}
\end{equation}

where $\alpha_h$ is the halation intensity ($0$ to $1$, default $0.18$) and
$\beta_c$ the channel bias. Because $\int h = 1$, \eqref{eq:haladd} preserves
total energy exactly when $S = E_c$, and the attenuation term correctly darkens
the core of an extremely bright highlight as its energy is redistributed
outward. Omitting the attenuation term --- as an additive glow does --- inflates
total scene energy and makes strong halation settings brighten the entire image,
a defect that manifests as "the halation slider is also an exposure slider."

\begin{figure}[htbp]
\centering
\begin{tikzpicture}
\begin{axis}[
  width=\columnwidth, height=4.6cm,
  xlabel={radius $r$ ($\mu$m at film plane)},
  ylabel={normalized $h(r)$},
  xmin=0, xmax=400, ymin=0, ymax=1.05,
  grid=both, grid style={draw=emLine!20},
  legend style={font=\scriptsize, at={(0.98,0.98)}, anchor=north east, draw=emLine!40},
  label style={font=\scriptsize}, tick label style={font=\scriptsize},
  samples=250,
]
\addplot[emAcc, very thick, domain=0:400] {exp(-x/95)};
\addlegendentry{$\ell_R = 95\,\mu$m}
\addplot[emGreen, very thick, domain=0:400] {exp(-x/59)};
\addlegendentry{$\ell_G = 59\,\mu$m}
\addplot[emBlue, very thick, domain=0:400] {exp(-x/42)};
\addlegendentry{$\ell_B = 42\,\mu$m}
\addplot[emLine, thick, dashed, domain=0:400] {exp(-(x*x)/(2*95*95))};
\addlegendentry{Gaussian, same $\sigma$}
\end{axis}
\end{tikzpicture}
\caption{Halation point spread functions. The dashed Gaussian falls to negligible values by $250\,\mu$m; the exponential retains a visible skirt far beyond. That skirt is the visual signature of halation and is why a Gaussian bloom does not read as film.}
\label{fig:halpsf}
\end{figure}

\subsection{Real-Time Evaluation}
\label{sec:recursive-blur}

A direct convolution with $h_\ell$ at $\ell = 95\,\mu$m on a 48\,MP render
corresponds to a kernel radius of roughly $150$ pixels, which is
$\sim\!70{,}000$ taps per pixel. This is not viable. Three strategies were
evaluated.

\paragraph{FFT convolution} Exact and $O(N \log N)$, but requires two complex
transforms of a 48\,MP image, which exceeds the memory budget of NFR-03 and is
awkward to tile.

\paragraph{Recursive IIR filtering} A separable first-order recursive filter
realizes $e^{-|x|/\ell}$ exactly in one dimension at $O(1)$ per pixel. Applied
separably it produces $e^{-(|x|+|y|)/\ell}$, which is not isotropic --- it has
visible diamond-shaped anisotropy that is objectionable on point highlights.
Deriche's third-order approximation \cite{deriche1993,young1995} improves
isotropy but not enough at these radii.

\paragraph{Fitted mip pyramid (chosen)} The isotropic exponential is
approximated as a weighted sum of Gaussians,

\begin{equation}
h_\ell(r) \;\approx\; \sum_{j=1}^{J} w_j\, \mathcal{N}\!\left(r;\, s_j \ell\right),
\qquad \sum_j w_j = 1,
\label{eq:gaussmix}
\end{equation}

each of which is separable, isotropic, and --- critically --- can be evaluated
at a mip level chosen so that its kernel is only a few pixels wide. With $J = 4$
levels spaced by factors of four in $\sigma$, the frequency-domain fit of
\eqref{eq:gaussmix} to \eqref{eq:halotf} has a maximum relative error below
$3\%$ over $[0, 1/(2\ell)]$, which is far below AC-4 tolerance.

The coefficients $(w_j, s_j)$ are produced by an offline least-squares fit in the
frequency domain, stored in the stock profile, and regenerated by the profile
authoring tool whenever the target PSF changes. Representative values for the
$J = 4$ fit are given in Table~\ref{tab:gaussmix}.

\begin{table}[htbp]
\caption{Representative Gaussian Mixture Fit to \eqref{eq:halotf}}
\label{tab:gaussmix}
\begin{center}
\renewcommand{\arraystretch}{1.15}
\begin{tabular}{@{}lccc@{}}
\toprule
\textbf{Term $j$} & $s_j$ ($\times\ell$) & $w_j$ & \textbf{Mip level} \\
\midrule
1 & $0.42$ & $0.451$ & 0 \\
2 & $1.15$ & $0.318$ & 1 \\
3 & $3.30$ & $0.163$ & 3 \\
4 & $9.60$ & $0.068$ & 5 \\
\bottomrule
\multicolumn{4}{@{}p{0.92\columnwidth}@{}}{\footnotesize Produced by the offline frequency-domain fitter; regenerate rather than hand-edit. Mip level is chosen so each Gaussian has $\sigma \le 3$ pixels at that level.}
\end{tabular}
\end{center}
\end{table}

The resulting pass structure is: build a five-level pyramid of the source term
$S$ by successive $2\times$ box-plus-tent downsamples; apply a small separable
Gaussian at each retained level; upsample with bilinear filtering and accumulate
with weights $w_j$; combine per \eqref{eq:haladd}. Total cost is approximately
$1.33\times$ the base-resolution pass count due to the geometric series, and the
memory overhead is one-third of a full-resolution single-channel texture.

The ring component, when $\omega > 0$, is evaluated on the level-3 pyramid tap
using a small annular kernel of 16 taps arranged on a Poisson-disc ring, since
the ring is spatially large but the mip level makes it only a few pixels wide.

\subsection{Parameter Summary}

\begin{table}[htbp]
\caption{Halation Parameters}
\label{tab:halparams}
\begin{center}
\renewcommand{\arraystretch}{1.15}
\begin{tabularx}{\columnwidth}{@{}ll>{\raggedright\arraybackslash}X@{}}
\toprule
\textbf{Symbol} & \textbf{Default} & \textbf{Meaning} \\
\midrule
$\tau$        & $1.6$   & Threshold, relative exposure \\
$\kappa_h$    & $0.15$  & Threshold knee softness \\
$q$           & $1.0$   & Source exponent \\
$\ell_R$      & $95\,\mu$m & Red scattering length \\
$\ell_G/\ell_R$ & $0.62$ & Green ratio \\
$\ell_B/\ell_R$ & $0.44$ & Blue ratio \\
$\alpha_h$    & $0.18$  & Intensity \\
$\boldsymbol{\beta}$ & $(1.0, 0.42, 0.22)$ & Channel bias \\
$\omega$      & $0.05$  & Base-reflection ring weight \\
$t_b$         & $125\,\mu$m & Base thickness (drives $r_{\min}$) \\
\bottomrule
\end{tabularx}
\end{center}
\end{table}

The user-facing controls are threshold, radius (a multiplier on $\ell_R$ with the
ratios held fixed), intensity, and color bias (a single warm--cool parameter that
moves $\boldsymbol{\beta}$ along a locus rather than exposing three channels).
The remaining parameters live in the stock profile.
